% Ausgeschriebene Hilfsbeweise für die beiden Pfropfungsbedingungen.
\FormulaThmDeltaK[Ein angehängtes Adressbit erzeugt nicht die Wurzel]{%
  p\in\BinaryAddresses\dsep i\in\{0,1\}\vdash
    \WordConcat p{\LetterWord i}\notin\{\EmptyWord{\{0,1\}}\}
}{B27BinaryAppendNotRoot}{\DeltaRow{Mengen}{p\dsep i}}
\begin{tabproofwide}
  \proofstepwidestar[1]{p\in\BinaryAddresses}{\rA}
  \proofstepwidestar[2]{i\in\{0,1\}}{\rA}
  \proofstepwidestar[1]{p\in\FiniteWords{\{0,1\}}}{\FormulaRefAuto{BinaryAddressWordSetDef}[def]{1}}
  \proofstepwidestar[1,2]{\sigma_{\{0,1\}}(p,i)\ne\EmptyWord{\{0,1\}}}{\FormulaRefAuto{EarlyWordSuccessorNonempty}[thm]{3,2}}
  \proofstepwidestar[1,2]{\sigma_{\{0,1\}}(p,i)=\WordConcat p{\LetterWord i}}{\FormulaRefAuto{WordModelSuccessorConcatenation}[thm]{3,2}}
  \proofstepwidestar[1,2]{\WordConcat p{\LetterWord i}\ne\EmptyWord{\{0,1\}}}{\rIE{5,4}}
  \proofstepwidestar[1,2]{\WordConcat p{\LetterWord i}\notin\{\EmptyWord{\{0,1\}}\}}{\FormulaRefAuto{x \notin \{a\} \eqvdash x \neq a}[thm]{6}}
\end{tabproofwide}

\FormulaThmDeltaK[Ein vorangestelltes Adressbit erzeugt nicht die Wurzel]{%
  p\in\BinaryAddresses\dsep i\in\{0,1\}\vdash
    \AddressPrefix i p\ne\EmptyWord{\{0,1\}}
}{B27BinaryPrefixNotRoot}{\DeltaRow{Mengen}{p\dsep i}}
\begin{tabproofwide}
  \proofstepwidestar[1]{p\in\BinaryAddresses}{\rA}
  \proofstepwidestar[2]{i\in\{0,1\}}{\rA}
  \proofstepwidestar[1]{p\in\FiniteWords{\{0,1\}}}{\FormulaRefAuto{BinaryAddressWordSetDef}[def]{1}}
  \proofstepwidestar[2]{\LetterWord i\in\NonemptyWords{\{0,1\}}}{\FormulaRefAuto{LetterWordNonempty}[thm]{2}}
  \proofstepwidestar[2]{\LetterWord i\in\FiniteWords{\{0,1\}}}{\FormulaRefAuto{NonemptyWordFinite}[thm]{4}}
  \proofstepwidestar[2]{\WordLength{\LetterWord i}\in\mathbb N}{\FormulaRefAuto{FiniteWordLengthNatural}[thm]{5}}
  \proofstepwidestar[1]{\WordLength p\in\mathbb N}{\FormulaRefAuto{FiniteWordLengthNatural}[thm]{3}}
  \proofstepwidestar[2]{\WordLength{\LetterWord i}\ne0}{\rAEb{\FormulaRefAuto{NonemptyWordMembership}[thm]{4}}}
  \proofstepwidestar[1,2]{\WordLength{\LetterWord i}+\WordLength p\ne0}{\FormulaRefAuto{PeanoAddNonzeroLeft}[thm]{6,7,8}}
  \proofstepwidestar[1,2]{\WordConcat{\LetterWord i}p\in\FiniteWords{\{0,1\}}}{\FormulaRefAuto{WordConcatenationClosure}[thm]{5,3}}
  \proofstepwidestar[1,2]{\WordLength{\WordConcat{\LetterWord i}p}=\WordLength{\LetterWord i}+\WordLength p}{\FormulaRefAuto{WordConcatenationLength}[thm]{5,3}}
  \proofstepwidestar[1,2]{\WordLength{\WordConcat{\LetterWord i}p}\ne0}{\rIE{\FormulaRefAuto{x=y\vdash y=x}[thm]{11},9}}
  \proofstepwidestar[1,2]{\WordConcat{\LetterWord i}p\ne\EmptyWord{\{0,1\}}}{\rRE{12,\rLREb{\FormulaRefAuto{FiniteWordLengthNonzeroIffNonempty}[thm]{10}}}}
  \proofstepwidestar[1,2]{\AddressPrefix i p=\WordConcat{\LetterWord i}p}{\FormulaRefAuto{BinaryAddressPrefixDef}[def]{2,1}}
  \proofstepwidestar[1,2]{\AddressPrefix i p\ne\EmptyWord{\{0,1\}}}{\rIE{\FormulaRefAuto{x=y\vdash y=x}[thm]{14},13}}
\end{tabproofwide}

\FormulaThmDeltaK[Die beiden Präfixbilder in einer Pfropfung]{%
  \begin{aligned}[t]
    &X,Y\in\powerset(\BinaryAddresses)\dsep p\in\BinaryAddresses\vdash{}\\[-2pt]
    &(\AddressPrefix0p\in\AddressGraft X Y\leftrightarrow p\in X)\land{}\\[-2pt]
    &(\AddressPrefix1p\in\AddressGraft X Y\leftrightarrow p\in Y)
  \end{aligned}
}{B27GraftPrefixMembership}{\DeltaRow{Mengen}{X\dsep Y\dsep p\dsep q}}
\begin{proof}
\emergencystretch=2em
Nach \FormulaRefAuto{x\in\powerset(A)\eqvdash x\subseteq A}[thm]
liegen \(X,Y\) im Träger der Präfixabbildungen; deren Typisierung
liefert \FormulaRefAuto{BinaryAddressPrefixDef}[def].
Sei \(\AddressPrefix0p\in\AddressGraft X Y\). Die Definition
\FormulaRefAuto{BinaryAddressGraftDef}[def] und das Vereinigungskriterium
\FormulaRefAuto{x\in A\cup B\eqvdash x\in A\lor x\in B}[thm]
geben drei Fälle. Die Zugehörigkeit zur Einermenge der Wurzel ist
wegen \FormulaRefAuto{B27BinaryPrefixNotRoot}[thm] und
\FormulaRefAuto{x\in\{a\}\eqvdash x=a}[thm] ausgeschlossen.
Im zweiten Fall liefert
\FormulaRefAuto{C\subseteq A\dsep F\colon A\to B\dsep y\in F[C]\vdash\exists x\in C\;y=F(x)}[thm]
einen Zeugen \(q\in X\) mit \(\AddressPrefix0p=\AddressPrefix0q\).
\FormulaRefAuto{BinaryAddressPrefixWordInjective}[thm] gibt \(p=q\),
also \(p\in X\). Im dritten Fall gäbe es \(q\in Y\) mit
\(\AddressPrefix0p=\AddressPrefix1q\).
\FormulaRefAuto{BinaryAddressPrefixLetterInjective}[thm] ergäbe
\(0=1\), entgegen \FormulaRefAuto{PeanoZeroNeOne}[thm].

Ist umgekehrt \(p\in X\), liefert derselbe Bildmengensatz in
Einführungsrichtung
\FormulaRefAuto{C\subseteq A\dsep F\colon A\to B\dsep\exists x\in C\;y=F(x)\vdash y\in F[C]}[thm]
mit dem Zeugen \(p\) die Zugehörigkeit
\(\AddressPrefix0p\in\AddressPrefixMap0[X]\).
Zweimalige Vereinigungseinführung und
\FormulaRefAuto{BinaryAddressGraftDef}[def] geben
\(\AddressPrefix0p\in\AddressGraft X Y\).
Für das zweite Präfixbild werden \(0,X\) und \(1,Y\) vertauscht:
Der Wurzelfall bleibt ausgeschlossen, der gleich markierte Bildfall
liefert \(p\in Y\), und der fremd markierte Fall ergäbe \(1=0\),
das durch Symmetrie der Gleichheit wieder \(0=1\) liefert.
Die beiden bewiesenen Richtungen werden jeweils durch \rLRI{} und
die beiden Äquivalenzen durch \rAI{} verbunden.
\end{proof}

\FormulaThmDeltaK[Präfixabschluss der gepfropften Adressmenge]{%
  \begin{aligned}[t]
    &\GoodAddressSet X\dsep\GoodAddressSet Y\vdash{}\\[-2pt]
    &\forall q\in\AddressGraft X Y\setminus\{\EmptyWord{\{0,1\}}\}\,
      \exists p\in\AddressGraft X Y\,\exists i\in\{0,1\}\;q=\WordConcat p{\LetterWord i}
  \end{aligned}
}{B27GraftPrefixClosure}{\DeltaRow{Mengen}{X\dsep Y\dsep q\dsep p\dsep i\dsep u\dsep v\dsep j}}
\begin{proof}
\emergencystretch=2em
\FormulaRefAuto{GoodBinaryAddressSetCarrierAxiom}[ax] liefert die
Trägerzugehörigkeiten von \(X,Y\).
Sei \(q\in\AddressGraft X Y\setminus\{\EmptyWord{\{0,1\}}\}\).
\FormulaRefAuto{BinaryAddressGraftDef}[def], das Vereinigungskriterium
\FormulaRefAuto{x\in A\cup B\eqvdash x\in A\lor x\in B}[thm]
und das Differenzkriterium
\FormulaRefAuto{x\in A\setminus B\vdash x\notin B}[thm]
geben \(q\in\AddressPrefixMap0[X]\) oder
\(q\in\AddressPrefixMap1[Y]\).

Im ersten Fall liefert das Bildmengenkriterium
\FormulaRefAuto{C\subseteq A\dsep F\colon A\to B\dsep y\in F[C]\vdash\exists x\in C\;y=F(x)}[thm]
einen Zeugen \(u\in X\) mit \(q=\AddressPrefix0u\).
Falls \(u=\EmptyWord{\{0,1\}}\), ist
\(q=\LetterWord0=\WordConcat{\EmptyWord{\{0,1\}}}{\LetterWord0}\)
nach \FormulaRefAuto{BinaryAddressPrefixDef}[def] und den beiden
Neutralitätsgesetzen \FormulaRefAuto{WordConcatenationIdentity}[thm].
Die Wurzel liegt in der Pfropfung, da sie in deren erster
Vereinigungskomponente liegt. Die gesuchten Zeugen sind also
\(p=\EmptyWord{\{0,1\}}\), \(i=0\).

Falls \(u\ne\EmptyWord{\{0,1\}}\), liegt \(u\) in
\(X\setminus\{\EmptyWord{\{0,1\}}\}\), nach
\FormulaRefAuto{x \notin \{a\} \eqvdash x \neq a}[thm] und
\FormulaRefAuto{x\in A\setminus B\eqvdash x\in A\land x\notin B}[thm].
\FormulaRefAuto{GoodBinaryAddressSetPrefixAxiom}[ax] liefert
\(v\in X\) und \(j\in\{0,1\}\) mit
\(u=\WordConcat v{\LetterWord j}\). Setze
\(p=\AddressPrefix0v\); nach
\FormulaRefAuto{B27GraftPrefixMembership}[thm] liegt \(p\) in der
Pfropfung. \FormulaRefAuto{BinaryAddressPrefixDef}[def] und
\FormulaRefAuto{WordConcatenationAssociative}[thm] ergeben
\[
 q=\WordConcat{\LetterWord0}{\WordConcat v{\LetterWord j}}
  =\WordConcat{\WordConcat{\LetterWord0}v}{\LetterWord j}
  =\WordConcat p{\LetterWord j}.
\]
Die Typisierungen folgen aus der Trägerzugehörigkeit von \(v\),
\FormulaRefAuto{LetterWordTyping}[thm] sowie
\FormulaRefAuto{WordConcatenationClosure}[thm].
Somit sind \(p\) und \(j\) die gesuchten Zeugen.
Der zweite Bildfall verwendet denselben Nachweis mit \(Y,1\):
Für \(u=\EmptyWord{\{0,1\}}\) sind die Zeugen Wurzel und \(1\),
sonst \(p=\AddressPrefix1v\), \(j\).
\rEE{} entlässt die gewählten Bild- und Vorgängerzeugen;
\rOE{} verbindet die Fälle, und \rUI{} schließt
über das beliebige \(q\).
\end{proof}

\FormulaThmDeltaK[Volle Verzweigung der gepfropften Adressmenge]{%
  \begin{aligned}[t]
    &\GoodAddressSet X\dsep\GoodAddressSet Y\vdash{}\\[-2pt]
    &\forall p\in\AddressGraft X Y\,
      (\WordConcat p{\LetterWord0}\in\AddressGraft X Y
        \leftrightarrow\WordConcat p{\LetterWord1}\in\AddressGraft X Y)
  \end{aligned}
}{B27GraftFullness}{\DeltaRow{Mengen}{X\dsep Y\dsep p\dsep u\dsep j}}
\begin{proof}
\emergencystretch=2em
Sei \(p\in\AddressGraft X Y\). Wie im vorigen Beweis liefern
\FormulaRefAuto{BinaryAddressGraftDef}[def], das Vereinigungs- und
das Einermengenkriterium die drei Fälle Wurzel, linkes Präfixbild und
rechtes Präfixbild. Die nötigen Trägerzugehörigkeiten stammen aus
\FormulaRefAuto{GoodBinaryAddressSetCarrierAxiom}[ax].

Für \(p=\EmptyWord{\{0,1\}}\) liegen nach
\FormulaRefAuto{GoodBinaryAddressSetRootAxiom}[ax] die Wurzeln in
\(X\) und \(Y\). \FormulaRefAuto{B27GraftPrefixMembership}[thm]
zeigt daher \(\AddressPrefix0{\EmptyWord{\{0,1\}}}\) und
\(\AddressPrefix1{\EmptyWord{\{0,1\}}}\) in der Pfropfung.
\FormulaRefAuto{WordConcatenationIdentity}[thm] identifiziert diese
Wörter mit \(\WordConcat p{\LetterWord0}\) bzw.
\(\WordConcat p{\LetterWord1}\). Beide Aussagen sind wahr, also
folgen beide Implikationen und ihre Äquivalenz.

Im linken Bildfall liefert der Bildmengensatz
\FormulaRefAuto{C\subseteq A\dsep F\colon A\to B\dsep y\in F[C]\vdash\exists x\in C\;y=F(x)}[thm]
ein \(u\in X\) mit \(p=\AddressPrefix0u\). Für jedes
\(j\in\{0,1\}\) liefern \FormulaRefAuto{BinaryAddressPrefixDef}[def]
und \FormulaRefAuto{WordConcatenationAssociative}[thm]
\[
 \WordConcat p{\LetterWord j}
 =\AddressPrefix0{\WordConcat u{\LetterWord j}}.
\]
Die zugehörigen Worttypisierungen folgen aus
\FormulaRefAuto{LetterWordTyping}[thm] und
\FormulaRefAuto{WordConcatenationClosure}[thm]. Daher gilt nach
\FormulaRefAuto{B27GraftPrefixMembership}[thm]
\[
 \WordConcat p{\LetterWord j}\in\AddressGraft X Y
 \leftrightarrow\WordConcat u{\LetterWord j}\in X.
\]
Für \(j=0\) und \(j=1\) verbindet
\FormulaRefAuto{GoodBinaryAddressSetFullnessAxiom}[ax] die rechten
Seiten. Äquivalenzersetzung \rLRS{} liefert die verlangte
Äquivalenz der linken Seiten. Im rechten Bildfall wird der Bildzeuge
\(u\in Y\) mit \(p=\AddressPrefix1u\) gewählt; dieselbe Rechnung
führt auf das Verzweigungsaxiom von \(Y\).
Die Bildzeugen werden durch \rEE{} entlassen, die drei Fälle durch
\rOE{} verbunden und das beliebige \(p\) durch \rUI{} gebunden.
\end{proof}
